% ED26-SRC-036-019: Turner 2008, complete selected chapter 759--776.
\chapter{Nonmonotonic Causal Theories and Action Histories}
\label{sec:p16-nonmonotonic-causal-theories}

A constraint on possible states does not necessarily explain an action's
effects. Turner's 2008 account distinguishes a fact's being true from its
being caused under a description. A causal rule supplies a sufficient
condition for its head to be caused; an explanatory model must be uniquely
determined by all the heads activated in that world
\cite[pp.~759--776]{Turner2008NonmonotonicCausalLogic}.
The result is a semantics of fully explained worlds or histories, not an
observational probability law or an identification of what caused what.
\index{causal theory!nonmonotonic}\index{causal explanation}
\index{frame problem}\index{action!ramification}

This treatment gives the source's finite-domain theory, replacement and
completion results, all six action cases, displayed C+ descriptions, and
default, higher-order and universal-causation comparisons. All formal results
promised here have full reciprocal annex proofs. Broader first-order completion,
full C+ grammar and general modal context-equivalence referrals retain their
actual cited boundaries.

\section{Values, activated explanations and models}
\label{sec:p16-nct-semantics}

\begin{hbdefinition}{def:p16-nct-signature}{Finite-domain propositional language}
A signature \(\sigma\) is a set of constants, each with a nonempty finite
domain \(\operatorname{Dom}(c)\) of value symbols. Atoms are \(c=v\), with
\(v\in\operatorname{Dom}(c)\); formulas are Boolean combinations of atoms.
An interpretation \(I\) assigns one domain value to every constant.
It satisfies \(c=v\) iff \(I(c)=v\). The signature need not be finite.
The empty signature has its one empty interpretation.

A Boolean constant has domain \(\{t,f\}\); \(c\) abbreviates \(c=t\),
and \(\neg c\) is equivalent to \(c=f\).
Locally the source's \(\equiv\) is object-formula biconditional,
\(\not\equiv\) its negation, and \(\supset\) material implication, not
elementary equivalence of structures.
Let \(\operatorname{Mod}(\Gamma)\) be the interpretations satisfying every
formula of \(\Gamma\), and \(\Gamma\models F\) the usual all-model consequence.
\end{hbdefinition}

\begin{hbdefinition}{def:p16-nct-theory}{Causal rules and explanatory models}
A causal rule is \(\phi\Leftarrow\psi\), with head \(\phi\) and body
\(\psi\). Its sign constructs a rule, not a material implication.
For a set \(T\) of such rules define
\[
 T^I=\{\phi:\phi\Leftarrow\psi\in T,\ I\models\psi\}.
\]
The activation test uses the body. An interpretation is an explanatory
model exactly when \(\operatorname{Mod}(T^I)=\{I\}\).
The heads must both hold and determine the world uniquely.
Consistency, or truth of each corresponding material implication, is weaker.
\end{hbdefinition}

This is the selector printed in the source's opening on p.~759 and used in
its cases. The isolated p.~763 display tests the head instead; the active
definition keeps the consistent body test. The cycling domain below is
\(\{0,1,2\}\), as the source's rules and p.~770 C+ declaration require,
rather than the incompatible domain printed above those rules on p.~767.
The retained source bytes are unchanged.

\begin{hbtheorem}{thm:p16-nct-basic}{proved=annex:p16.nct-basic}{Truth, explanation and nonmonotonicity}
\(I\) is explanatory exactly when, for every formula \(F\),
\(I\models F\) iff \(T^I\models F\).
Every explanatory model satisfies every \(\psi\supset\phi\) for
\(\phi\Leftarrow\psi\in T\); the converse fails.
Adding a causal rule can add explanatory models.
\end{hbtheorem}

For the source's three-valued \(p\) with domain \(\{1,2,3\}\),
\(\{p=1\Leftarrow p=1\}\) explains only value \(1\).
Adding \(p=2\Leftarrow p=2\) adds the world \(p=2\); neither rule
explains \(3\). This expands possible worlds under added explanations.

\begin{hbtheorem}{thm:p16-nct-extensions}{proved=annex:p16.nct-extensions}{Constraints and fresh definitions}
Adding \(\bot\Leftarrow F\) removes exactly the explanatory models satisfying
\(F\), and adds none. It does not explain \(\neg F\).

For fresh Boolean \(d\) and old-signature \(F\), adding
\(d\equiv F\Leftarrow\top\) gives each old model its unique extension
with \(d=t\) iff \(F\) holds, and gives no others.

For fresh \(d\) with distinct values \(v_1,\ldots,v_r\), \(r\geq1\),
and pairwise incompatible old formulas \(F_2,\ldots,F_r\), add
\[
 d=v_1\Leftarrow d=v_1,\qquad d=v_i\Leftarrow F_i\ (2\leq i\leq r).
\]
Each old model extends uniquely with \(v_i\) when \(F_i\) holds and
otherwise \(v_1\). Dropping incompatibility can make the extension impossible.
\end{hbtheorem}

These are the source's constraint, definitional-extension and overridden
default-value constructions \cite[pp.~763--765]{Turner2008NonmonotonicCausalLogic}.

\section{Strong equivalence and replacement}
\label{sec:p16-nct-strong-equivalence}

\begin{hbdefinition}{def:p16-nct-se}{Contexts and SE-models}
Equivalence means equality of explanatory-model sets.
Strong equivalence means that \(S\cup R\) and \(T\cup R\) are equivalent
for every causal theory \(R\) on the same signature.
An SE-model of \(T\) is a pair \((I,J)\) with
\(I\models T^I\) and \(J\models T^I\).
The first world activates the heads; both worlds test those same heads,
not \(T^J\). Neither must itself be an explanatory model.
\end{hbdefinition}

\begin{hbtheorem}{thm:p16-nct-strong}{proved=annex:p16.nct-strong}{SE-model characterization}
Two causal theories are strongly equivalent iff their SE-models are identical.
This holds for set-sized signatures and theories; the annex constructs the
separating context explicitly. Replacing a head and a body by respectively
equivalent formulas preserves strong equivalence.
\end{hbtheorem}

Ordinary equivalence is insufficient. The source's
\(S=\{p\Leftarrow q,q\Leftarrow p\}\) and
\(T=\{p\Leftarrow p,q\Leftarrow q\}\) each explain only \(p=q=t\).
With \(R=\{\neg p\Leftarrow\neg p\}\), \(T\cup R\) additionally explains
\(\neg p,q\), whereas \(S\cup R\) does not
\cite[pp.~765--766]{Turner2008NonmonotonicCausalLogic}.
The causal pair test is not automatically the strong-equivalence semantics
of a logic program.

\section{Completion and finite computation}
\label{sec:p16-nct-completion}

\begin{hbdefinition}{def:p16-nct-completion}{Definite theories and completion}
A theory is definite when every head is an atom or \(\bot\) and each atom
heads only finitely many rules. Boolean negative heads can be expressed as
false-value atoms. An atom is trivial when its constant has a singleton
domain. Completion contains \(\neg F\) for each constraint
\(\bot\Leftarrow F\), and, for every nontrivial atom \(A\),
\[
                  A\equiv\bigvee_{A\Leftarrow B\in T}B.
\]
An empty disjunction is \(\bot\). The head-finiteness assumption makes each
formula finite even when the theory or signature is infinite.
\end{hbdefinition}

\begin{hbtheorem}{thm:p16-nct-completion}{proved=annex:p16.nct-completion}{Exact completion}
The explanatory models of a definite causal theory are exactly its
completion models. No acyclicity or tightness assumption is needed.
Trivial atoms must be excluded: the empty theory on the singleton domain
\(\{0\}\) has its unique world, while \(c=0\equiv\bot\) is unsatisfiable.
\end{hbtheorem}

For Boolean \(p\Leftarrow p\), completion includes \(p\equiv p\)
\emph{and} \(\neg p\equiv\bot\), so only the true world is explained.
The normal program \(p\leftarrow p\) instead has empty least reduct model
under \cref{def:p09-stable-model}; its Clark completion is just
\(p\leftrightarrow p\), allowing both values. Causal completion, program
completion and stable models are different tests
\cite[pp.~766--768]{Turner2008NonmonotonicCausalLogic}.

\begin{hbmethod}{def:p16-nct-sat-method}{Finite value encoding}
For finite explicitly listed domains and rules, form completion and give
every value atom a Boolean variable. Require at least one value and pairwise
at most one value for each constant; translate the Boolean connectives.
These domain constraints and completion represent exactly the original
assignments. Subsequent definitional clause encoding requires its own fresh
variables and model-projection argument; it does not license dropping
negative-value completions or empty head lists.
\end{hbmethod}

\begin{hbtheorem}{thm:p16-nct-complexity}{proved=annex:p16.nct-complexity}{Finite decision classifications}
For well-formed finite theories with formulas and domains explicitly encoded,
definite explanatory-model existence is \(\mathsf{NP}\)-complete,
strong equivalence is \(\mathsf{coNP}\)-complete, and ordinary equivalence
of arbitrary causal theories is \(\Pi_2^p\)-complete, under polynomial
many-one reductions. Hardness already holds on Boolean signatures.
The benchmark imports are
\cref{thm:p04-cook-levin,thm:p04-fixed-qbf-completeness}.
No complexity claim for an unencoded infinite theory is made.
\end{hbtheorem}

These are Turner's stated classifications; the annex supplies explicit
certificates and reductions. SAT compilation is not proof that a rule base
is an adequate physical model.

\section{All six action cases}
\label{sec:p16-nct-action-cases}

Time subscripts are distinct constants, not temporal operators.
For initial values and action values, exogeneity adds
\(c_0=v\Leftarrow c_0=v\) and \(a_t=w\Leftarrow a_t=w\).
For an inertial fluent, at each specified successor time add
\[
                    c_{t+1}=v\Leftarrow c_{t+1}=v\land c_t=v
\]
for every declared value. Exogeneity stipulates available explanations;
it does not investigate the actual outside causes.

\subsection{Concurrent handles and soup}
Use Boolean \(upL_0,upR_0,sp_0,upL_1,upR_1,sp_1,liftL_0,liftR_0\).
Include exogeneity for all initial fluents/actions and inertia for all three
fluents, then all four laws
\[
\begin{gathered}
 upL_1\Leftarrow liftL_0,\quad upR_1\Leftarrow liftR_0,\\
 sp_0\Leftarrow upL_0\not\equiv upR_0,\quad
 sp_1\Leftarrow upL_1\not\equiv upR_1.
\end{gathered}
\]
A material state constraint is equivalent to its contrapositive but supplies
no explanation for a handle rising. Initially both handles are down and the
soup unspilled. Lifting only the left spills the soup; lifting both does not.
The alternative right handle rising without its lift has no support.
This is the source's fully specified bowl theory
\cite[pp.~760--764]{Turner2008NonmonotonicCausalLogic}.

\subsection{A three-value cycle}
Let \(p_0,p_1\) have domain \(\{0,1,2\}\) and \(a_0\) be Boolean.
Include initial/action exogeneity and inertia, then
\[
 p_1=1\Leftarrow p_0=0\land a_0,\quad
 p_1=2\Leftarrow p_0=1\land a_0,\quad
 p_1=0\Leftarrow p_0=2\land a_0.
\]
The non-tautological completion equations are
\(p_1=v\equiv(p_1=v\land p_0=v)\lor(p_0=v^{-}\land a_0)\),
where \(v^{-}\) precedes \(v\) in \(0,1,2,0\).
An action advances the value; no action preserves it
\cite[pp.~767,770]{Turner2008NonmonotonicCausalLogic}.

\subsection{Coin tossing}
Declare \(coin_0,coin_1\) with domain \(\{heads,tails\}\), Boolean \(toss_0\),
initial/action exogeneity, inertia, and
\[
              coin_1=v\Leftarrow coin_1=v\land toss_0
                              \quad(v\in\{heads,tails\}).
\]
A toss gives either successor an explanation. It does not assign a
probability to either outcome \cite[p.~768]{Turner2008NonmonotonicCausalLogic}.

\subsection{Putting an object}
For \(k>1\), \(loc_0,loc_1\) range over \(\{1,\ldots,k\}\);
all \(put(v)_0\) are Boolean. Add initial/action exogeneity, location inertia,
and \(loc_1=v\Leftarrow put(v)_0\) for every \(v\).
Two different put actions activate incompatible values. Their joint
impossibility is implied, not imposed by pairwise prohibition rules
\cite[pp.~768--769]{Turner2008NonmonotonicCausalLogic}.

\subsection{Dominos}
For \(k\geq1\), use Boolean \(tip_0\) and \(up(d)_t\),
\(1\leq d\leq k,\ 0\leq t\leq k\).
Include \(up(d)_0\Leftarrow\top\), tip exogeneity and both Boolean
inertia rules for \(0\leq t<k\). Add
\[
\begin{gathered}
 \neg up(1)_1\Leftarrow tip_0,\\
 \neg up(d+1)_{t+2}\Leftarrow\neg up(d)_{t+1}\land up(d)_t
       \quad(1\leq d<k,\ 0\leq t\leq k-2).
\end{gathered}
\]
A fall, rather than another action, triggers the next domino at the next
time. The final family is empty at \(k=1\).
The finite synchronous horizon is part of the stipulated model
\cite[p.~769]{Turner2008NonmonotonicCausalLogic}.

\subsection{A pendulum's tendency}
On horizon \(0,\ldots,n\), \(n\geq0\), use Boolean \(left_t\),
\(0\leq t\leq n\), and \(hold_t\), \(0\leq t<n\).
Include exogeneity of \(left_0\) and all holds;
for every \(0\leq t<n\) include exactly
\[
\begin{array}{ll}
 left_{t+1}\Leftarrow left_{t+1}\land\neg left_t,&
 \neg left_{t+1}\Leftarrow\neg left_{t+1}\land left_t,\\
 left_{t+1}\Leftarrow hold_t\land left_t,&
 \neg left_{t+1}\Leftarrow hold_t\land\neg left_t.
\end{array}
\]
The first row supports change; the second forces persistence when held.
Adding ordinary inertia would alter the theory
\cite[p.~770]{Turner2008NonmonotonicCausalLogic}.

\begin{hbtheorem}{thm:p16-nct-action-models}{proved=annex:p16.nct-action-models}{Exact action outcomes and counts}
At the all-false initial bowl state, actions neither, left only, right only,
both have respectively successors
\((upL_1,upR_1,sp_1)=(f,f,f),(t,f,t),(f,t,t),(t,t,f)\), uniquely.
The cycle has six histories: three unchanged and three advanced.
The coin has six: two untossed unchanged, four tossed.
Location has \(k(k+1)\): each initial location with zero or one put action.
Dominos have exactly two: no tip and all remain up; or a tip and
\(up(d)_t\) iff \(d>t\).
The pendulum preserves its value exactly when held, otherwise flips;
each initial value/hold sequence has one history, giving \(2^{n+1}\).
\end{hbtheorem}

\section{The displayed C+ fragment}
\label{sec:p16-nct-cplus}

The source defines states by models of \(T(D,0)\), transitions by models
of \(T(D,1)\), and length-\(n\) paths by models of \(T(D,n)\).
Its cycling description declares \texttt{inertial p}, \texttt{exogenous a},
and \texttt{a causes p=1 if p=0}, \texttt{a causes p=2 if p=1},
\texttt{a causes p=0 if p=2}. Its bowl description declares inertial
\(upL,upR,sp\), exogenous \(liftL,liftR\), both direct lifting effects,
and \texttt{caused sp if upL} \(\not\equiv\) \texttt{upR}
\cite[pp.~770--771]{Turner2008NonmonotonicCausalLogic}.

\begin{hbdefinition}{def:p16-nct-cplus-fragment}{Complete translation of those descriptions}
Use fluent constants \(c_t\) at \(0\leq t\leq n\), and action constants
\(a_t\) at \(0\leq t<n\). Add initial/action exogeneity and all declared
inertial rules.
Translate \texttt{a causes F if G} as
\(F_{t+1}\Leftarrow a_t\land G_t\), and \texttt{caused F if G}
as \(F_t\Leftarrow G_t\) at every state time.
Subscript substitution applies to all fluent occurrences together.
A missing condition is \(\top\); the displayed \(F\) is a value atom
or Boolean fluent. This specifies the two displayed descriptions, not
the whole C+ grammar or its later extensions.
\end{hbdefinition}

\begin{hbtheorem}{thm:p16-nct-cplus-paths}{proved=annex:p16.nct-cplus-paths}{States, transitions and paths}
For these descriptions, explanatory models of \(T(D,n)\) are exactly
the paths of length \(n\) in the indicated transition system, retaining
all action values on edge labels.
The cycle has unchanged and advancing edges from each state.
Bowl states satisfy \((upL\not\equiv upR)\supset sp\); each of four action
choices has its unique successor. Horizon zero has states and no actions.
\end{hbtheorem}

This gives no planning algorithm, optimization objective or stochastic law.

\section{Default correspondence and its restriction}
\label{sec:p16-nct-default}

For Boolean signatures translate \(\phi\Leftarrow\psi\) to the
prerequisite-free default \(\frac{:\psi}{\phi}\), with empty background facts.
For a world \(I\), \(\operatorname{Th}(I)\) is all classical formulas
true there: a consistent complete theory.

\begin{hbtheorem}{thm:p16-nct-default}{proved=annex:p16.nct-default}{Complete consistent default extensions}
\(I\) is explanatory iff \(\operatorname{Th}(I)\) is a Reiter extension of
the translated defaults. Incomplete extensions need not be worlds.
For \(\{p\Leftarrow q,q\Leftarrow q,\neg q\Leftarrow\neg q\}\), the defaults
have exactly \(\operatorname{Th}(\{p,q\})\) and
\(\operatorname{Th}(\{\neg q\})\); only the first is complete, corresponding
to the only explanatory world.
\end{hbtheorem}

The proof uses \cref{def:p08-default-theory}, not a circular test of
nonempty prerequisites in a candidate.
The source connects these inertia rules to Reiter's 1980 discussion,
distinguishing later successful action encodings from the different
Yale Shooting encoding \cite[pp.~771--772]{Turner2008NonmonotonicCausalLogic}.
This attributed history is not a theorem about every default action theory.

\section{Higher-order and universal causation}
\label{sec:p16-nct-representations}

\begin{hbdefinition}{def:p16-nct-higher-order}{Explainable-symbol alternatives}
In a classical signature designate a finite list \(N\) of explainable
nonlogical constants, with other symbols fixed.
A finite nonpropositional theory has rules
\(\phi_i(N,\bar x_i)\Leftarrow\psi_i(N,\bar x_i)\), where \(\bar x_i\)
lists all free variables. For fresh variables \(m\) of matching sorts/arities,
put
\[
 T^*(m)=\bigwedge_i\forall\bar x_i
             (\psi_i(N,\bar x_i)\supset\phi_i(m,\bar x_i)).
\]
Only heads use the alternative assignment. Over fixed object domains and
fixed nonexplainable symbols, full typed semantics ranges alternatives over
all corresponding objects/functions/predicates. Equality \(m=N\) is
componentwise extensional equality. The model condition is
\[
                         \forall m\,(T^*(m)\equiv m=N).
\]
\end{hbdefinition}

\begin{hbtheorem}{thm:p16-nct-higher-order}{proved=annex:p16.nct-higher-order}{Unique assignments and Boolean overlap}
This formula says exactly that the actual explainable-symbol assignment
uniquely satisfies every head activated at it, with all free-variable
instances covered.
If all nonlogical symbols are explainable Boolean propositional constants,
it agrees with explanatory-model semantics.
For finite Boolean theories, enumerating all \(2^{|N|}\) assignments for
\(m\) eliminates the Boolean quantifiers, possibly increasing the size
exponentially by this construction.
\end{hbtheorem}

Turner's referral to Lifschitz's broader nonpropositional completion
does not specify its full definite-theory hypotheses here; it is not
an unrestricted first-order theorem inferred from propositional completion
\cite[pp.~772--773]{Turner2008NonmonotonicCausalLogic}.

\begin{hbdefinition}{def:p16-nct-ucl}{Universal causation in S5 structures}
Add formulas \(\mathsf C F\), read as \(F\) being caused.
An S5 structure is \((I,\mathcal S)\) with \(I\in\mathcal S\), all worlds
of one signature. Atoms use \(I\), connectives their truth tables, and
\[
 (I,\mathcal S)\models\mathsf C F
   \quad\Longleftrightarrow\quad
 (J,\mathcal S)\models F\text{ for all }J\in\mathcal S.
\]
Nested formulas use this recursion. An \(I\)-model of a UCL theory satisfies
it at that structure. \(I\) is a causal model when \((I,\{I\})\) is its
unique \(I\)-model. This nonmonotonic selection on S5 structures is
not an additional monotonic S5 axiom or a definition of knowledge.
\end{hbdefinition}

\begin{hbtheorem}{thm:p16-nct-ucl}{proved=annex:p16.nct-ucl}{The exact modal embedding}
For any finite-domain causal theory \(T\), its explanatory models equal
the causal models of
\[
                       \{\psi\supset\mathsf C\phi:
                                         \phi\Leftarrow\psi\in T\}.
\]
A finite-signature restriction is unnecessary.
\end{hbtheorem}

For general UCL, the source describes SE-structures satisfying the theory
both at \((I,\mathcal S)\) and at \((I,\{I\})\), referring to the separate
context-equivalence development. A set-of-worlds test for nested modal
formulas is not the causal-rule pair test silently reused
\cite[pp.~773--774]{Turner2008NonmonotonicCausalLogic}.

\section{Consultation and causal boundaries}
\label{sec:p16-nct-use}

\begin{hbmethod}{def:p16-nct-use}{Select and test the actual explanatory question}
Specify domains, horizon, exogenous initial/action facts, inertia or tendencies,
direct/indirect laws and constraints.
Activate bodies in a candidate history, then check both truth and uniqueness
of the heads. For finite definite theories use completion with domain
constraints. For safe replacement use SE-models, not isolated model sets.

In \cref{def:p16-finite-scm}, structural functions assign endogenous variables
and a disturbance law supports interventions by equation replacement.
The coin theory here admits two successors for one initial/action state
without a probability law or an undeclared disturbance. Adding disturbances
can give a structural model, but does not identify these semantics.
Likewise \cref{def:p09-stable-model,def:p09-clark-completion} use least
reduct models or predicate definitions rather than causal uniqueness.
The Reiter correspondence above selects only complete consistent
prerequisite-free extensions. No bridge makes all causation monotonic,
deductive, statistical or interchangeable.
\end{hbmethod}

The source's strategic notion describes conditions under which facts count
as caused, often stipulating outside causes for initial facts and actions.
It does not directly reason about what causes what or infer a causal
mechanism from observational data.
The full captured original C+ grammar and its state/history constructions
are developed in \cref{sec:p16-cplus-action-language}, separately from
this chapter's two displayed descriptions. Elaboration tolerance, additive
fluents, probabilistic/game extensions and broader nonpropositional results
are not consequences of the six stipulated cases.
